\documentclass[11pt,a4paper]{article}
\usepackage{ftm}

\title{Navier–Stokes\\[4pt]
       \large Caderno de Lacunas Formais — Programa FTM}
\author{Federated Trivalent Manifold Research Programme}
\date{2026}

\begin{document}
\maketitle
\tableofcontents
\bigskip

%% ---------------------------------------------------------------
\section{Enunciado oficial}
%% ---------------------------------------------------------------

\begin{definition}[Problema de Navier–Stokes]
Para velocidade incompressível $u:\mathbb R^3\times[0,\infty)\to\mathbb R^3$
e pressão $p:\mathbb R^3\times[0,\infty)\to\mathbb R$, as equações são
\[
  \partial_t u+(u\cdot\nabla)u = -\nabla p+\nu\Delta u+f,
  \qquad \nabla\cdot u=0,
\]
com dado inicial $u_0\in\mathcal S(\mathbb R^3)$ e força $f\in C^\infty$.
O problema pergunta se existe sempre uma solução suave e globalmente
limitada, ou se pode ocorrer breakdown em tempo finito.
\end{definition}

\bigskip

%% ---------------------------------------------------------------
\section{Programas de pesquisa e encaixe FTM}
%% ---------------------------------------------------------------

\begin{programs}

\item \textbf{Regularidade condicional — critérios de Prodi–Serrin.}
Se $u\in L^p(0,T;L^q(\mathbb R^3))$ com $2/p+3/q=1$, $q\ge3$,
então $u$ é regular em $[0,T]$.
Escauriaza–Seregin–Šverák (2003) provaram que blow-up no critério
crítico $L^3$ é impossível, i.e.\ $\|u(\cdot,t)\|_{L^3}\to\infty$
ao se aproximar de um tempo de blow-up.

\begin{gap}
Os critérios de Prodi–Serrin são condicional: eles excluem
blow-up abaixo do limiar crítico, mas não provam regularidade global
sem controle prévio de $u$ em norma crítica. \OPEN
\end{gap}

\item \textbf{Integração convexa e não-unicidade.}
Buckmaster–Vicol (2019) construíram, via integração convexa e
convex $h$-principle, soluções de Leray fracas não únicas
em $L^2_t L^3_x$ abaixo da regularidade de Onsager.
O análogo é a conjunção Onsager: soluções Hölder $C^\alpha$
conservam energia sse $\alpha>1/3$ (provado em ambas as direções).

\begin{gap}
A não-unicidade de soluções fracas não implica breakdown da solução
forte; o gap entre regularidade de Leray e regularidade Prodi–Serrin
crítica permanece aberto. \OPEN
\end{gap}

\item \textbf{Blow-up numérico rigoroso — cenário Hou–Luo.}
Hou–Luo (2014) identificam, via simulação numérica, um cenário
de blow-up potencial em domínio cilíndrico com fronteira.
O programa de verificação intervalar (computação aritmética rigorosa
com arredondamento controlado) busca certificar esse cenário.

\begin{gap}
Nenhuma verificação rigorosa (intervalar ou outra) prova que
o cenário Hou–Luo produz blow-up em tempo finito para dado suave.
Simulações não substituem prova de existência de singularidade. \OPEN
\end{gap}

\item \textbf{Espaços críticos — Koch–Tataru.}
Koch–Tataru (2001) provaram bem-postura local em $BMO^{-1}$
para dados pequenos; a solução existe e é suave enquanto a norma
$BMO^{-1}$ permanece controlada.
O espaço $BMO^{-1}$ é crítico pela escala de Navier–Stokes.

\end{programs}

%% ---------------------------------------------------------------
\section{Discretização FTM}
%% ---------------------------------------------------------------

\begin{definition}[Família FTM refinável para Navier–Stokes]
Uma família refinável $\{F_h\}_{h\to0}$ com espaços $V_h,Q_h$
e operadores discretos
\[
  \nabla_h,\quad \operatorname{div}_h,\quad \Delta_h,\quad P_h
\]
deve satisfazer: compatibilidade de colagem em interseções,
um complexo discreto exato (de Rham discreto) e
a condição inf-sup de Babuška–Brezzi.
\end{definition}

%% ---------------------------------------------------------------
\section{Obrigações formais FTM para Navier–Stokes}
%% ---------------------------------------------------------------

\begin{obligation}[NS1 — Consistência]\label{ob:NS1}
Para $u$ solução suave,
$\|\text{resíduo}_h(u)\|_{V_h^*}\to0$ quando $h\to0$,
com taxa quantificada.
\end{obligation}

\begin{obligation}[NS2 — Estabilidade energética]\label{ob:NS2}
\[
  \|u_h(t)\|_2^2 + 2\nu\int_0^t\|\nabla_h u_h\|_2^2\,ds
  \;\le\; \|u_h(0)\|_2^2 + \int_0^t(f,u_h)\,ds.
\]
\end{obligation}

\begin{obligation}[NS3 — Compacidade e limite]\label{ob:NS3}
De qualquer sequência $h_k\to0$, extrair subsequência tal que
$u_{h_k}\to u$ forte o suficiente para passar ao limite nas
equações; o limite $u$ satisfaz Navier–Stokes no sentido
de distribuições.
\end{obligation}

\begin{obligation}[NS4 — Estimativa crítica]\label{ob:NS4}
Obter controle uniforme em norma de escala crítica
(e.g.\ $L^3$, $BMO^{-1}$) suficiente para provar regularidade global,
ou construir sequência $u_{h_k}$ que converge para uma solução fraca
satisfazendo exatamente as hipóteses do enunciado oficial e
que desenvolve singularidade em tempo finito.
\end{obligation}

\begin{remark}
Convergência para solução de Leray fraca (NS1–NS3) não resolve
a pergunta de regularidade.
Simulação de vórtices não prova blow-up.
NS4 é a lacuna decisiva e está \OPEN.
\end{remark}

\MATURITY{M1 conceitual. NS1–NS4 todas ABERTAS; NS4 é a lacuna central.}

\end{document}
